\documentclass[11pt,a4paper]{article}

\def\courseTitle{Industrial Security (Bachelor)}
\def\sectionNumber{10}
\def\sectionTitle{Incident Response and Recovery}
\def\sheetKind{Solution Sheet}

% =====================================================================
% Shared preamble for exercise sheets and solution sheets.
%
% Define BEFORE \input{}-ing this file:
%   \def\courseTitle{...}
%   \def\sectionNumber{...}
%   \def\sectionTitle{...}
%   \def\sheetKind{Exercise Sheet}   or  Solution Sheet
%
% Optional:
%   \def\courseAuthor{...}
%
% The caller must issue \documentclass{...} (typically article) BEFORE
% loading this preamble.
% =====================================================================

\usepackage[utf8]{inputenc}
\usepackage[T1]{fontenc}
\usepackage[english]{babel}
\usepackage[a4paper, margin=2.2cm]{geometry}
\usepackage{graphicx}
\usepackage{listings}
\usepackage{xcolor}
\usepackage{tikz}
\usepackage{booktabs}
\usepackage{array}
\usepackage{enumitem}
\usepackage{titlesec}
\usepackage{fancyhdr}
\usepackage{hyperref}
\usepackage{amsmath}
\usepackage{amssymb}
\usepackage{tcolorbox}
\tcbuselibrary{skins, breakable}
\usepackage{needspace}

% Discourage solo lines split across pages.
\widowpenalty=10000
\clubpenalty=10000
\raggedbottom

% =====================================================================
% Shared TikZ styles for the Industrial Security lecture series.
% Loaded by every slide deck and exercise sheet that uses TikZ.
% =====================================================================

\usetikzlibrary{
  arrows.meta,
  shapes,
  shapes.geometric,
  shapes.symbols,
  shapes.multipart,
  positioning,
  fit,
  calc,
  backgrounds,
  decorations.pathmorphing,
  decorations.pathreplacing,
  decorations.text,
  patterns,
  matrix,
  shadows
}

% --- colour palette --------------------------------------------------
\definecolor{isOT}{HTML}{1F4E79}     % OT (deep blue)
\definecolor{isIT}{HTML}{6F2DA8}     % IT (purple)
\definecolor{isSafety}{HTML}{C00000} % safety red
\definecolor{isNet}{HTML}{2E7D32}    % network green
\definecolor{isAttack}{HTML}{B23A48} % attack red
\definecolor{isAccent}{HTML}{E08E0B} % amber
\definecolor{isMute}{HTML}{6B6B6B}   % grey
\definecolor{isLight}{HTML}{EDF2F8}  % light background
\definecolor{isPLC}{HTML}{2C5282}
\definecolor{isHMI}{HTML}{2F855A}
\definecolor{isSCADA}{HTML}{6B46C1}

% --- generic node styles --------------------------------------------
\tikzset{
  isnode/.style={
    draw, rounded corners=2pt, minimum height=8mm, minimum width=18mm,
    font=\small, align=center, thick
  },
  isbox/.style={isnode, fill=isLight},
  plc/.style={isnode, fill=isPLC!15, draw=isPLC, thick},
  hmi/.style={isnode, fill=isHMI!15, draw=isHMI, thick},
  scada/.style={isnode, fill=isSCADA!15, draw=isSCADA, thick},
  sensor/.style={isnode, fill=isNet!10, draw=isNet, minimum height=6mm, minimum width=14mm, font=\scriptsize},
  actuator/.style={isnode, fill=isAccent!15, draw=isAccent, minimum height=6mm, minimum width=14mm, font=\scriptsize},
  firewall/.style={isnode, fill=isAttack!15, draw=isAttack, thick},
  server/.style={isnode, fill=isMute!15, draw=isMute!80, thick},
  switch/.style={isnode, fill=isLight, draw=black!60, minimum height=6mm, minimum width=14mm, font=\scriptsize},
  workstation/.style={isnode, fill=white, draw=isOT, thick},
  attacker/.style={isnode, fill=isAttack!20, draw=isAttack, thick, font=\small\bfseries},
  inetnode/.style={draw, ellipse, minimum width=24mm, minimum height=10mm, font=\scriptsize, fill=isLight, thick},
  process/.style={isnode, fill=isOT!15, draw=isOT, thick, minimum width=24mm},
  zone/.style={
    draw, rounded corners=4pt, thick, dashed, inner sep=4mm
  },
  conduit/.style={-{Stealth[length=2.5mm]}, thick, isNet!80!black},
  attack/.style={-{Stealth[length=2.5mm]}, thick, isAttack, dashed},
  flow/.style={-{Stealth[length=2.5mm]}, thick, isOT!70!black},
  netline/.style={thick, black!60},
  axisline/.style={-{Stealth[length=2mm]}, thick, black!70},
  tag/.style={font=\scriptsize\itshape, fill=white, inner sep=1pt},
  level/.style={
    draw, fill=isLight, minimum width=0.85\linewidth, minimum height=8mm,
    font=\small, anchor=west
  },
  brace/.style={decorate, decoration={brace, amplitude=4pt, raise=2pt}},
  legendbox/.style={draw, rounded corners=2pt, fill=isLight, inner sep=2mm, font=\scriptsize, align=left}
}

% --- handy macros ----------------------------------------------------
% \plcicon, \hmiicon etc as simple labelled boxes; useful inline.
\newcommand{\plcicon}[1][PLC]{\tikz[baseline=-0.6ex]{\node[plc, minimum width=10mm, minimum height=5mm, font=\scriptsize, inner sep=1pt]{#1};}}
\newcommand{\hmiicon}[1][HMI]{\tikz[baseline=-0.6ex]{\node[hmi, minimum width=10mm, minimum height=5mm, font=\scriptsize, inner sep=1pt]{#1};}}
\newcommand{\fwicon}[1][FW]{\tikz[baseline=-0.6ex]{\node[firewall, minimum width=10mm, minimum height=5mm, font=\scriptsize, inner sep=1pt]{#1};}}


\providecommand{\courseAuthor}{Industrial Security Lecture Series}
\providecommand{\courseInstitute}{Department of Computer Science}
\providecommand{\companyLogo}{logo}

% --- Corporate branding overrides ------------------------------------
% Picks up logo / author / brand colours from the repo-root branding.tex
% when present; falls back to the defaults above otherwise.
\IfFileExists{../../../branding.tex}{% =====================================================================
% Corporate / institutional branding -- single file to customise the
% look of the entire course (slides, notes, exercises, labs).
%
% HOW TO REBRAND IN UNDER A MINUTE:
%   1. Replace `branding/logo.pdf` with your own logo
%      (PDF preferred; PNG and JPG also work -- name it `logo.<ext>`).
%   2. (Optional) change the two colours below to your brand palette.
%   3. (Optional) override the author/institute lines below.
%   4. Re-run `make` at the repo root. That's it.
%
% Everything in this file has sensible defaults; you can delete a line
% to fall back to the built-in defaults, or comment it out.
%
% This file is loaded automatically by the slides, exercise, and lab
% preambles -- you never need to \input it manually.
% =====================================================================

% --- Logo --------------------------------------------------------------
% Path is resolved from each leaf-build directory; the default points at
% the shared `branding/` folder at the repo root. If you put your logo
% somewhere else, set the full or relative path here (without extension):
\renewcommand{\companyLogo}{../../../branding/logo}

% --- Author / institute (appears on the title page footer) ------------
\renewcommand{\courseAuthor}{}
\renewcommand{\courseInstitute}{}

% --- Brand colours ----------------------------------------------------
% slideAccent      = primary brand colour (title bands, accent rule)
% slideSecondary   = secondary brand colour (section badge, highlight)
% Both use HTML hex codes. Keep enough contrast against white text.
\definecolor{slideAccent}{HTML}{00205B}      % deep navy
\definecolor{slideSecondary}{HTML}{00A0DC}   % light blue accent
}{}

% --- listings -------------------------------------------------------
\definecolor{codebg}{rgb}{0.96,0.96,0.96}
\lstset{
  backgroundcolor=\color{codebg},
  basicstyle=\ttfamily\footnotesize,
  breaklines=true,
  frame=single,
  numbers=left,
  numberstyle=\tiny\color{black!50},
  showstringspaces=false,
  tabsize=2
}

% --- header / footer ------------------------------------------------
\pagestyle{fancy}
\fancyhf{}
\renewcommand{\headrulewidth}{0.4pt}
\renewcommand{\footrulewidth}{0pt}
\lhead{\small \courseTitle}
\rhead{\small Section \sectionNumber{} -- \sheetKind}
\cfoot{\thepage}

% --- task / solution boxes -----------------------------------------
% `before={...}` guards every task/solution box with \needspace so the
% box doesn't start with only one or two lines of breathing room at the
% bottom of a page. If less than ~9 lines remain, LaTeX bumps the box
% to the next page; otherwise it starts here and breaks normally if it
% runs long.
% Boxes are `breakable` so very long solutions don't blow the page,
% but `before={\needspace{14\baselineskip}}` pushes them to a fresh
% page whenever fewer than ~14 lines remain. That covers the vast
% majority of single-question splits in practice.
\newtcolorbox{taskbox}[1]{
  enhanced, breakable,
  colback=isLight, colframe=isOT,
  fonttitle=\bfseries\color{white},
  coltitle=white,
  colbacktitle=isOT,
  title={#1},
  arc=2pt, boxrule=0.6pt,
  left=2mm, right=2mm, top=2mm, bottom=2mm,
  before={\par\vspace{2mm}\needspace{14\baselineskip}\par},
  after={\par\vspace{2mm}}
}

\newtcolorbox{solutionbox}{
  enhanced, breakable,
  colback=isNet!5, colframe=isNet,
  fonttitle=\bfseries\color{white},
  coltitle=white,
  colbacktitle=isNet,
  title={Solution},
  arc=2pt, boxrule=0.6pt,
  left=2mm, right=2mm, top=2mm, bottom=2mm,
  before={\par\vspace{2mm}\needspace{14\baselineskip}\par},
  after={\par\vspace{2mm}}
}

\newtcolorbox{hintbox}{
  enhanced, breakable,
  colback=isAccent!10, colframe=isAccent,
  fonttitle=\bfseries\color{white}, coltitle=white, colbacktitle=isAccent,
  title={Hint},
  arc=2pt, boxrule=0.6pt
}

\newtcolorbox{warnbox}{
  enhanced, breakable,
  colback=isAttack!8, colframe=isAttack,
  fonttitle=\bfseries\color{white}, coltitle=white, colbacktitle=isAttack,
  title={Caution},
  arc=2pt, boxrule=0.6pt
}

% --- section formatting --------------------------------------------
\titleformat{\section}{\Large\bfseries\color{isOT}}{\thesection}{0.5em}{}
\titleformat{\subsection}{\large\bfseries\color{isOT!80!black}}{\thesubsection}{0.5em}{}

% --- title block macro ----------------------------------------------
\newcommand{\sheetheader}{%
  \begin{center}
    {\Large\bfseries \courseTitle}\\[0.3em]
    {\large \sheetKind{} -- Section \sectionNumber}\\[0.2em]
    {\large \sectionTitle}\\[0.2em]
    \rule{\linewidth}{0.4pt}
  \end{center}
  \vspace{0.5em}
}


\begin{document}
\sheetheader

\noindent
\textit{How to use this sheet:} the solutions are sample answers, not the only
correct ones. Each solution box ends with a \textbf{Rubric} hint listing the
points that must appear for full marks. If the student's answer covers the
rubric with different but defensible choices, mark it as correct.

\begin{solutionbox}
\textbf{Exercise 10.1 -- Historian ransomware tabletop.}
\textit{T+0--30 min, Detection \& Analysis:} operator confirms HMI still
controls process (safety unaffected); shift supervisor opens incident ticket;
on-call OT security pages incident commander. \textit{T+30--90 min,
Containment:} IC isolates Historian VLAN at L3 firewall but leaves SCADA-HMI
path intact so operators keep visibility; PCAP and ransom note hashed and
preserved; backup admin pulls last-night immutable backup off the shelf and
does \emph{not} restore yet. \textit{T+90--180 min, Escalation:} plant manager
declares Severity~1; legal and PR briefed; national CSIRT receives
NIS2 Art.\,23(4)(a) early warning placeholder; cyber-insurance hotline opened.
\textit{T+180--240 min, Decision:} plant manager (not security) decides whether
to continue production blind on trends or step down to a safe state; comms
sends internal stand-up note; 24-hour early-warning text drafted for sign-off
before the legal deadline. \textbf{Rubric:} (1) safety check first, (2) named
decision owner per step, (3) preservation before remediation, (4) NIS2 24-hour
clock acknowledged, (5) plant manager owns the production decision.
\end{solutionbox}

\begin{solutionbox}
\textbf{Exercise 10.2 -- Forensic scope.}
\textit{PLC (M580 / S7-1500):} \emph{can} pull running project via vendor tool,
diagnostic buffer, firmware version, CPU mode; \emph{cannot} pull a forensic
memory image (no documented dump interface, vendor often refuses), no per-byte
disk image; \emph{compensating}: signed baseline comparison, upstream PCAP of
S7/Modbus writes, engineering audit trail. \textit{HMI:} can image disk and
RAM (Windows host), event logs, USB history; cannot trust the HMI clock alone;
compensating: SIEM correlation, AD logon events. \textit{Engineering WS:} can
image disk, RAM, TIA / RSLogix audit trail, project-file timestamps, USB
history, browser history; cannot recover deleted ladder-logic versions if
project-file history is off; compensating: vendor cloud sync, git history,
PLC-side ``last download'' record. \textit{Network:} full PCAP at the iDMZ and
L2 access switch (if SPAN/TAP exists), NetFlow, IDS alerts, firewall syslog,
DHCP/DNS logs; cannot get what was never captured; compensating: historical
NetFlow, neighbour-switch ARP/MAC tables. \textbf{Rubric:} each row has all
three columns; mention the PLC memory-dump gap (NIST SP 800-86 ``volatile
data'' is largely infeasible on controllers); mention chain of custody.
\end{solutionbox}

\begin{solutionbox}
\textbf{Exercise 10.3 -- RACI for ``Shut down Unit 2''.}
Sample assignment: \emph{Responsible} -- SCADA / DCS engineer (executes the
controlled shutdown); \emph{Accountable} -- plant manager (owns production
risk and the safety case); \emph{Consulted} -- security operations, HSE
officer, legal; \emph{Informed} -- PR / communications. Security is not
accountable because the cost of a wrong shutdown (lost batch, restart hazards,
flare events) typically dwarfs the cyber risk, and only the plant manager has
the statutory authority for an operating-permit-relevant decision. HSE is
consulted (not accountable) because their role is to veto on safety grounds,
not to own the production call. \textbf{Rubric:} exactly one Accountable;
Accountable is the plant manager; HSE has a veto path; PR is Informed, not
Consulted, in the decision moment.
\end{solutionbox}

\begin{solutionbox}
\textbf{Exercise 10.4 -- Recovery validation of a re-flashed M580.}
(1)~Firmware hash matches the vendor-published value (tool: vendor checksum
utility; pass: byte-identical). (2)~Project upload from the PLC equals the
signed baseline project (tool: Unity Pro / Control Expert compare; pass: zero
diff). (3)~Diagnostic buffer is clean and time-synced to the plant NTP source
(tool: vendor diag; pass: no unexplained entries since boot). (4)~Network
behaviour matches baseline -- only expected peers, expected ports (tool:
Zeek / Wireshark; pass: no new flows for 60 minutes idle). (5)~Loop test with
process simulator or on a bypassed loop confirms IO mapping and setpoints
unchanged (tool: control engineer + isolated test rig; pass: outputs within
tolerance). (6)~72-hour heightened-monitoring window with elevated SIEM rules
on the controller's VLAN (tool: SIEM; pass: zero high-severity alerts before
sign-off). \textbf{Rubric:} integrity check before behavioural check; project
compared to \emph{signed} baseline, not to ``what we had yesterday''; explicit
monitoring window after restart; plant manager (not security) signs the
production-ready statement.
\end{solutionbox}

\begin{solutionbox}
\textbf{Exercise 10.5 -- NIS2 early-warning sample text.}
\emph{``This is an early warning under Art.\,23(4)(a) of Directive (EU)
2022/2555. On 2026-05-12 at 22:15 CEST, [Utility] detected a confirmed
compromise of a SCADA jump host serving its drinking-water control network.
The incident is suspected to be caused by a malicious act; investigation is
ongoing and no attribution is asserted at this time. Based on current
information, cross-border impact is not expected, and drinking-water supply
to customers is not affected. A full incident notification under
Art.\,23(4)(b) will follow within 72 hours.''} \textbf{Rubric:} explicitly
labelled as Art.\,23(4)(a) early warning; states suspected malicious cause
without attribution; cross-border statement present; commits to the 72-hour
follow-up; 3--5 sentences; no speculation.
\end{solutionbox}

\begin{solutionbox}
\textbf{Exercise 10.6 -- RTO / RPO assignment.}
\begin{itemize}\setlength{\itemsep}{1pt}\small
  \item \textbf{Historian} -- RTO 24\,h, RPO 1\,h. Trend-gap tolerable; reporting continuity required.
  \item \textbf{MES} -- RTO 8\,h, RPO 4\,h. Order tracking drives throughput billing.
  \item \textbf{SCADA HMI} -- RTO 1\,h, RPO 0. Operators lose visibility of the live process.
  \item \textbf{SIS} -- RTO 0, RPO 0. Plant must trip rather than run without SIS.
  \item \textbf{Billing IT} -- RTO 48\,h, RPO 24\,h. Customers tolerate one-day delay.
  \item \textbf{Document mgmt} -- RTO 72\,h, RPO 24\,h. Change-mgmt only, not live ops.
\end{itemize}
\textbf{Rubric:} SIS has RTO/RPO of zero (plant trips, no degraded mode); HMI
RTO is the tightest live-ops system; billing and document management are the
loosest; every row carries a consequence anchor (safety / regulatory /
financial / reputational).
\end{solutionbox}

\begin{solutionbox}
\textbf{Exercise 10.7 -- Blameless vs.\ blameful opener.}
\textit{Blameless:} ``On 2026-04-18 at 09:42, a phishing email reached a
control-room workstation and was opened. The click triggered a download that
our endpoint controls did not block, and within 14 minutes the payload had
reached the Historian. This post-mortem focuses on the system conditions that
allowed a hostile message to reach a control-network host, the detection
controls that missed the lateral move, and the response steps that worked.
Our goal is to learn, not to assign blame: anyone in that seat could have
clicked, and the system must tolerate that.'' \textit{Blameful:} ``On
2026-04-18, operator J.\ Schmidt clicked a phishing email despite annual
training, causing the Historian outage. This document records the failure of
the operator to follow policy and recommends disciplinary action.''
\textit{Contrast:} the blameless version names the system gaps that need
fixing and keeps operators willing to self-report next time; the blameful
version drives reporting underground and leaves the email-gateway and
segmentation gaps untouched, so the same class of incident will recur.
\textbf{Rubric:} blameless opener names systems not people; blameful opener
names the individual; contrast paragraph identifies the reporting-chill
effect.
\end{solutionbox}

\begin{solutionbox}
\textbf{Exercise 10.8 -- Tabletop design for pre-positioning.}
\textit{Injects:} T+0 ``EDR alert: persistent C2 beacon from
ENG-WS-04''; T+15 ``DNS logs show beacon active for 47 days''; T+30 ``audit
trail: TIA Portal opened project file PROD\_REACTOR\_3 yesterday''; T+60
``threat-intel: beacon TLS fingerprint matches reported state actor''.
\textit{Discussion moments:} (1) Do we re-image now (and tip the adversary)
or watch (and accept risk)? (2) Do we file an NIS2 Art.\,23 early warning
when there is no impact yet but malicious cause is likely? (3) Do we bring
in the national CSIRT or handle internally? \textit{Success criteria:}
(a) team distinguishes ``observe'' from ``act'' deliberately, with named
owner; (b) chain-of-custody plan for the WS image is named before any
action; (c) NIS2 timing acknowledged and a draft early warning produced;
(d) plant-manager decision on Unit 3 continued operation is taken on
evidence, not feeling; (e) one concrete preparation gap is captured for the
lessons-learned register. \textbf{Rubric:} realistic timeline (not all in
one minute); at least one ``observe vs.\ act'' dilemma; explicit regulator
question; named scoring criteria.
\end{solutionbox}

\begin{solutionbox}
\textbf{Exercise 10.9 -- ``Pull the cable'' counter-examples.}
\textit{Scenario A (SIS):} a workstation that programs a Triconex / HIMA SIS
is suspected of compromise during a live batch. Pulling its network cable
strands the SIS in its current configuration with no engineering visibility,
and any subsequent process upset cannot be diagnosed or overridden.
\emph{Alternative:} freeze the workstation's user session, snapshot
volatile memory and disk while still connected, isolate at the SIS-zone
firewall after the batch reaches a safe hold point, and only then take the
host offline. \textit{Scenario B (Operator visibility):} the HMI server is
suspected of compromise mid-shift. Yanking the network blinds the operator
to alarms during a process transient; the resulting manual mis-handling can
be worse than the cyber incident itself. \emph{Alternative:} bring up the
warm-spare HMI on a clean VLAN, fail operators over, then quarantine the
suspect HMI for forensics. \textbf{Rubric:} at least one SIS / safety-loop
example; alternative is concrete (not just ``be careful''); alternative
preserves operator visibility and forensic evidence.
\end{solutionbox}

\begin{solutionbox}
\textbf{Exercise 10.10 -- Lessons-learned feedback artefacts.}
(1)~\textbf{Detection-rule updates} (owner: SOC lead; lives in SIEM / IDS
rule repo; trigger: rule merged and tagged with incident ID).
(2)~\textbf{Updated network diagram / asset inventory} (owner: OT
architect; lives in CMDB; trigger: diagram version bumped and approved by
change board). (3)~\textbf{Playbook revision} for the affected scenario
class (owner: IR lead; lives in IR repo / runbook tool; trigger: new
version published and team-read-receipted). (4)~\textbf{Tabletop scenario}
distilled from the real incident (owner: exercise coordinator; lives in
the exercise library; trigger: scenario scheduled into the next
quarter's tabletop). (5)~\textbf{Backup / recovery procedure update}
(owner: backup admin; lives in BCM runbook; trigger: dry-run restore
performed and signed off by plant manager). \textbf{Rubric:} each
artefact has owner / location / completion trigger; at least one feeds
detection, one feeds rehearsal, one feeds recovery; ``the post-mortem
report'' is not counted.
\end{solutionbox}

\end{document}
